%\documentclass[12pt,oneside]{amsart}
\documentclass{scrartcl}

%\documentclass{article}

\usepackage{amsthm,amsmath,amssymb}
\usepackage[numbers]{natbib}

%\usepackage[colorlinks]{hyperref}
\usepackage{hyperref}
\hypersetup{
    colorlinks,%
    citecolor=black,%
    filecolor=black,%
    linkcolor=black,%
    urlcolor=black
}
\usepackage{hypernat}
\usepackage[top=2.5cm, bottom=2.5cm, left=2.8cm,
right=2.8cm]{geometry}
\usepackage{graphicx}


%\setlength{\parskip}{0pt}
%\setlength{\parsep}{0pt}
%\setlength{\headsep}{0pt}
%\setlength{\topskip}{0pt}
%\setlength{\topmargin}{0pt}
%\setlength{\topsep}{0pt}
%\setlength{\partopsep}{0pt}

\linespread{1}

%\usepackage{marvosym}

%\textwidth = 6.5 in \textheight = 9.2 in \oddsidemargin = 0.0 in
%\evensidemargin = 0.0 in \topmargin = -0.0 in \headheight = 0.0 in
%\headsep = 0.0in \parskip = 0.0in \parindent = 0.0in
%\def\baselinestretch{0.871}


\newtheorem{theorem}{Theorem}[section]
\newtheorem{proposition}[theorem]{Proposition}
\newtheorem{problem}[theorem]{Problem}
\newtheorem{fact}[theorem]{Fact}
\newtheorem{lemma}[theorem]{Lemma}
\newtheorem{defn}[theorem]{Definition}
\newtheorem{corollary}[theorem]{Corollary}
\newtheorem{remark}{Remark}[section]
\newtheorem{example}[theorem]{Example}

\newcommand{\E}{{\mathbb E}}
\newcommand{\se}{{\mathcal{E}}}
\newcommand{\R}{{\mathbb R}}
\renewcommand{\P}{{\mathbb P}}
\newcommand{\ac}{{\mathcal{A}}}
\newcommand{\A}{{\mathcal{A}}}
\newcommand{\B}{{\mathcal{B}}}
\newcommand{\C}{{\mathcal{C}}}
\newcommand{\M}{{\mathcal{M}}}
\newcommand{\Mf}{{\mathfrak{M}}}
\newcommand{\T}{{\mathcal{T}}}
\newcommand{\Z}{{\mathcal{Z}}}
\newcommand{\K}{{\mathcal{K}}}
\newcommand{\lc}{{\mathcal{L}}}
\newcommand{\Ps}{{\mathcal{P}}}
\newcommand{\G}{{\mathcal{G}}}
\newcommand{\pa}{{\mathring{p}}}
\newcommand{\F}{{\cal F}}
\newcommand{\samp}{{\mathcal{X}}}
\newcommand{\X}{{\mathcal{X}}}
\newcommand{\hi}{{\hat{\Phi}}}
\newcommand{\data}{{\text{data}}}
\newcommand{\relu}{{\text{ReLU}}}

\newcommand{\ridge}{{\hat{\beta}_{\text{ridge}}(\lambda)}}
\newcommand{\lasso}{{\hat{\beta}_{\text{lasso}}(\lambda)}}
\newcommand{\ridgemu}{{\hat{\mu}_t^{\text{ridge}}(\lambda)}}
\newcommand{\lassomu}{{\hat{\mu}_t^{\text{lasso}}(\lambda)}}


\newcounter{rcnt}[section]
\renewcommand{\thercnt}{(\roman{rcnt})}

\def\qt#1{\qquad\text{#1}}

\def\argmin{\mathop{\rm argmin}}
\def\argmax{\mathop{\rm argmax}}


\setlength{\parskip}{1.4 \medskipamount}
%\sloppy
%\linespread{1.3}

\begin{document}

\title{STAT 238 - Bayesian Statistics  \\
  Lecture Thirty} 
\subtitle{Spring 2026, UC Berkeley}

\author{Aditya Guntuboyina}


\date{13 April 2026}

\maketitle

\section{The Gibbs Sampler for Gaussian Mixture Models}

We observe real-valued data $y_1, \dots, y_n$ and we consider the model:
\begin{align*}
  y_i \overset{\text{i.i.d}}{\sim} \sum_{j=1}^k w_j N(\mu_j, \sigma_j^2)
\end{align*}
with unknown parameters $(w_1, \dots, w_k)$ and $(\mu_j, \sigma_j^2)$
for $j = 1, \dots, k$.  $(w_1, \dots, w_k)$ is a probablity vector
(i.e., $w_j \geq 0$ and $\sum_j w_j = 1$).

We use the following priors:
\begin{align*}
  & (w_1, \dots, w_k) \sim \text{Dirichlet}(a_1, \dots, a_k) \\
  & \mu_1, \dots, \mu_k \overset{\text{i.i.d}}{\sim} N(m, s^2) \\
  & \sigma_1^2, \dots, \sigma_k^2 \overset{\text{i.i.d}}{\sim}
    IG(\alpha, \beta). 
\end{align*}
The default choices are $a_j = 1$ for all $j$, $m = 0$ and $s^2$ to be
large, and $\alpha, \beta$ to be near zero. Note that the density of
the Inverse Gamma distribution $IG(\alpha, \beta)$ is given by
$x^{-\alpha - 1} \exp(-\beta/x) I\{x > 0\}$. It is easy to check that
this is simply the density of $X^{-1}$ where $X \sim
\text{Gamma}(\alpha, \beta)$.  

To use the Gibbs sampler, we introduce latent variables $z_1, \dots,
z_n$ 
which represent group memberships. Specifically, $z_i$ represents
which of the $k$ Gaussians the observation $y_i$ is coming from. More
precisely, $z_i$ takes the values $1, \dots, k$ with probabilities
$w_1, \dots, w_k$ (i.e., $z_i \sim \text{Categorical}(w)$), and 
\begin{align*}
  y_i \mid z_i = j \sim N(\mu_j, \sigma_j^2). 
\end{align*}
The Gibbs sampler will then jointly sample from all the unknown
variables:
\begin{align*}
  w_1, \dots, w_k, \mu_1, \sigma_1^2, \dots, \mu_k, \sigma_k^2, z_1,
  \dots, z_n \mid y_1, \dots, y_n
\end{align*}
The full conditionals corresponding to all the variables can be
written in closed form as described below. We will use the notation $w
= (w_1, \dots, w_k)$, $\mu = (\mu_1, \dots, \mu_k)$ and $\sigma =
(\sigma_1, \dots, \sigma_k)$, also $y = (y_1, \dots, y_n)$ and $z =
(z_1, \dots, z_n)$. 

For the full conditional corresponding to $z_1, \dots, z_n$ (i.e., the
conditional distribution of $z_1, \dots, z_n$ given $y_1, \dots, y_n$
as well as $w_1, \dots, w_k, \mu_1, \sigma_1^2, \dots, \mu_k,
\sigma_k^2$) is given by:
\begin{align*}
  \P \left\{z_i = j \mid y_i, w, \mu, \sigma \right\} = r_{ij} := \frac{w_j
  \phi(y_i, \mu_j, \sigma_j^2)}{\sum_{j=1}^k w_j
  \phi(y_i, \mu_j, \sigma_j^2)}. 
\end{align*}
In other words,
\begin{align*}
  z_i \mid y_i, w, \mu, \sigma \sim \text{Categorical}(r_i) ~~
  \text{where $r_i = (r_{i1}, \dots, r_{ik})$}. 
\end{align*}
The full conditional of $w = (w_1, \dots, w_k)$ (i.e., the conditional
distribution of $w$ given $z, y, \mu, \sigma$) is:
\begin{align*}
  w \mid z, y, \mu, \sigma \sim \text{Dirichlet}(a_1 + n_1, \dots, a_k
  + n_k) ~~ \text{where $n_j := \sum_{i=1}^n I\{z_i = j\}$}. 
\end{align*}
The full conditional of $\mu_j$ (i.e., the conditional distribution of
$\mu_j$ given $z, y, \sigma$) is:
\begin{align*}
  \mu_j \mid z, y, \sigma \sim N \left(\frac{m/s^2 + \sum_{i : z_i =
  j} y_i/\sigma_j^2}{1/s^2 + n_j/\sigma_j^2}, \frac{1}{1/s^2 +
  n_j/\sigma_j^2} \right). 
\end{align*}
Note that we are also conditioning on $\sigma_j$ above. If we do not
condition on $\sigma_j$, then the conditional distribution of
$\mu_j$ will be $t$-distributed and not normal distributed. 

The full conditional of $\sigma_j^2$  (i.e., the conditional
distribution of $\sigma_j^2$ given $z, y, \mu$) is:
\begin{align*}
  \sigma_j^2 \mid z, y, \sigma \sim IG \left(\alpha + \frac{n_j}{2},
  \beta + \frac{1}{2} \sum_{i : z_i = j} (y_i - \mu_j)^2 \right), 
\end{align*}
where again $n_j = \sum_{i=1}^n I\{z_i = j\}$.

So the Gibbs sampler algorithm is given by:
\begin{enumerate}
\item Initialize the parameters $w^{(0)}, \mu^{(0)}, \sigma^{(0)}$.
\item Repeat the following for $t = 0, 1, 2, \dots$:
  \begin{enumerate}
  \item Generate $z_1^{(t)}, \dots, z_n^{(t)}$ via:
    \begin{align*}
      z_i^{(t)} \sim \text{Categorical} \left(\frac{w_j^{(t)}
  \phi \left(y_i, \mu^{(t)}_j, (\sigma_j^{(t)})^2 \right)}{\sum_{j=1}^k w_j^{(t)}
  \phi \left(y_i, \mu^{(t)}_j, (\sigma_j^{(t)})^2 \right)}, j = 1, \dots, k \right)
    \end{align*}
  \item Calculate
    \begin{align*}
      n_j^{(t)} := \sum_{i=1}^n I\{z_i^{(t)} = j\}  ~~ \text{ for $j =
      1, \dots, k$}.
    \end{align*}
  \item Generate $w^{(t+1)} = (w_1^{(t+1)}, \dots, w_k^{(t+1)})$ via
    \begin{align*}
      w^{(t+1)} \sim \text{Dirichlet}(a_1 + n_1^{(t)}, \dots, a_k +
      n_k^{(t)}). 
    \end{align*}
  \item Generate $\mu_1^{(t+1)}, \dots, \mu_k^{(t+1)}$ via
    \begin{align*}
      \mu_j^{(t+1)} \sim N \left(\frac{m/s^2 + \sum_{i : z^{(t)}_i =
  j} y_i/(\sigma_j^{(t)})^2}{1/s^2 + n_j^{(t)}/(\sigma_j^{(t)})^2},
      \frac{1}{1/s^2 + 
  n^{(t)}_j/(\sigma_j^{(t)})^2} \right). 
    \end{align*}
  \item Generate $\sigma_1^{(t+1)}, \dots, \sigma_k^{(t+1)}$ via
    \begin{align*}
   (\sigma_j^2)^{(t+1)} \sim IG \left(\alpha + \frac{n^{(t)}_j}{2},
  \beta + \frac{1}{2} \sum_{i : z^{(t)}_i = j} (y_i - \mu^{(t+1)}_j)^2 \right), 
    \end{align*}
  \end{enumerate}
\end{enumerate}

\section{The EM Algorithm}

The EM algorithm can be seen as a deterministic variant of the Gibbs
sampler. We shall look at the EM in more generality later. In the case
of fitting finite normal mixtures, the EM algorithm is given by the
following: 
\begin{enumerate}
\item Initialize the parameters $w^{(0)}, \mu^{(0)}, \sigma^{(0)}$.
\item Repeat the following for $t = 0, 1, 2, \dots$ until convergence:
  \begin{enumerate}
  \item \textbf{E-step.} Compute the responsibility of component $j$ for
  observation $i$:
    \begin{align*}
      r_{ij}^{(t)} := \frac{w_j^{(t)}
      \phi\left(y_i, \mu_j^{(t)}, (\sigma_j^{(t)})^2\right)}
      {\sum_{l=1}^k w_l^{(t)}
      \phi\left(y_i, \mu_l^{(t)}, (\sigma_l^{(t)})^2\right)},
      \quad j = 1, \dots, k.
    \end{align*}
  \item Compute the effective counts
    \begin{align*}
      n_j^{(t)} := \sum_{i=1}^n r_{ij}^{(t)}, \quad j = 1, \dots, k.
    \end{align*}
  \item \textbf{M-step.} Update the weights:
    \begin{align*}
      w_j^{(t+1)} := \frac{n_j^{(t)}}{n}, \quad j = 1, \dots, k.
    \end{align*}
  \item Update the means:
    \begin{align*}
      \mu_j^{(t+1)} := \frac{\sum_{i=1}^n r_{ij}^{(t)}\, y_i}{n_j^{(t)}},
      \quad j = 1, \dots, k.
    \end{align*}
  \item Update the variances:
    \begin{align*}
      (\sigma_j^2)^{(t+1)} := \frac{\sum_{i=1}^n r_{ij}^{(t)}
      \left(y_i - \mu_j^{(t+1)}\right)^2}{n_j^{(t)}},
      \quad j = 1, \dots, k.
    \end{align*}
  \end{enumerate}
\end{enumerate}

There are the following close connections between the EM algorithm and
the Gibbs sampler: 
\begin{itemize}
  \item The \textbf{E-step} replaces sampling $z_i \sim \text{Categorical}(\ldots)$
  with computing its expectation
  \begin{align*}
    r_{ij}^{(t)} = \mathbb{E}\left[I\{z_i = j\} \mid y_i, w^{(t)}, \mu^{(t)},
    \sigma^{(t)}\right],
  \end{align*}
  which are the same normalised weights used in the Categorical draw.

  \item The \textbf{effective count} $n_j^{(t)} = \sum_{i=1}^n r_{ij}^{(t)}$
  replaces the integer count $\sum_{i=1}^n I\{z_i^{(t)} = j\}$.

  \item The \textbf{weight update} replaces the Dirichlet draw and takes the
  maximum likelihood estimate
  \begin{align*}
    w_j^{(t+1)} = \frac{n_j^{(t)}}{n},
  \end{align*}
  which is equivalent to the posterior mean of the Dirichlet as the prior
  becomes uninformative.

  \item The \textbf{mean update} replaces the Gaussian conjugate draw and is
  the responsibility-weighted sample mean
  \begin{align*}
    \mu_j^{(t+1)} = \frac{\sum_{i=1}^n r_{ij}^{(t)}\, y_i}{n_j^{(t)}},
  \end{align*}
  the maximum likelihood estimate with no prior shrinkage toward $m$.

  \item The \textbf{variance update} replaces the inverse-Gamma draw and is
  the responsibility-weighted sample variance
  \begin{align*}
    (\sigma_j^2)^{(t+1)} = \frac{\sum_{i=1}^n r_{ij}^{(t)}
    \left(y_i - \mu_j^{(t+1)}\right)^2}{n_j^{(t)}},
  \end{align*}
  the maximum likelihood estimate with no $\alpha, \beta$ regularisation terms.
\end{itemize}



%\begin{thebibliography}{99}

%\bibitem[Roberts and Rosenthal(2004)]{RobertsRosenthal2004}
%Roberts, G. O. and Rosenthal, J. S. (2004).
%General state space Markov chains and MCMC algorithms.
%\textit{Probability Surveys}, \textbf{1}, 20--71.

%\bibitem[Meyn and Tweedie(2009)]{MeynTweedie2009}
%Meyn, S. P. and Tweedie, R. L. (2009).
%\textit{Markov Chains and Stochastic Stability}, 2nd ed.
%Cambridge University Press.


  
%\bibitem[Chib and Greenberg(1995)]{chib1995}
%Chib, S. and Greenberg, E. (1995).
%Understanding the Metropolis--Hastings algorithm.
%\textit{The American Statistician}, \textbf{49}, 327--335.


%\bibitem[MacKay and Peto(1995)]{MacKay1995HierarchicalDirichlet}
%D.~J.~C. MacKay and L.~C. Peto,
%\newblock ``A hierarchical Dirichlet language model,''
%\newblock \emph{Natural Language Engineering}, vol.~1,
%no.~3, pp.~289--308, 1995.

%\bibitem[Rubin(1981)]{Rubin1981BayesianBootstrap}
%D.~B. Rubin,
%\newblock ``The Bayesian bootstrap,''
%\newblock \emph{The Annals of Statistics}, vol.~9,
%no.~1, pp.~130--134, 1981.
  
%\bibitem[Fisher(1934)]{Fisher1934}
%R.~A. Fisher,
%\newblock ``Probability, likelihood and quantity of information in the logic of
%uncertain inference,''
%\newblock \emph{Proceedings of the Royal Society of London. Series~A,
%Containing Papers of a Mathematical and Physical Character}, vol.~146,
%no.~856, pp.~1--8, 1934.

%\bibitem[Efron(2010)]{Efron}
%Efron, B. (2010)
%\textit{Large-scale inference: empirical Bayes methods for estimation,
%testing, and prediction}. Cambridge University Press, 2010.

%\bibitem[Shumway and Stoffer(2017)]{ShumwayStoffer}
%Shumway, R. H. and Stoffer, D. S. (2017)
%\textit{Time Series Analysis and Its Applications: With R
%  Examples}. Springer, 4th edition. 
  
%     \bibitem[Jaynes(2003)]{Jaynes} Jaynes, E. T, (2003)
%  \textit{Probability theory: the logic of science}. Cambridge
%  University Press, 2003.
 
%  \bibitem[Sinai(1992)]{Sinai} Sinai, Y. G, (1992)
%  \textit{Probability theory: an introductory course}. Springer
%  Verlag, 1992.  

%\bibitem[{deGroot(1986)}]{DeGrootBlackwell}DeGroot, M. H. (1986)
%       A conversation with David Blackwell.
%\newblock \emph{Statistical Science}, \textbf{1}, 40--53.

%         \bibitem[Mosteller(1987)]{Mosteller} Mosteller, F, (1987)
%  \textit{Fifty challenging problems in probability with
%    solutions}. Courier Corporation, 1987.

%    \bibitem[MacKay(2003)]{MacKay} MacKay, D. J. C., (2003)
%  \textit{Information theory, inference, and learning
%  algorithms}. Cambridge University Press, 2003.

%\bibitem[{Lindley(2013)}]{Lindley}Lindley, D. V. (2013)
%   \textit{Understanding Uncertainty}. John Wiley and sons, 2013.


%\end{thebibliography}







\end{document}

%%% Local Variables:
%%% mode: latex
%%% TeX-master: t
%%% End:
